
\textbf{Exemple 1.5.} Soit @@M0@@ un schéma localement noethérien universellement caténaire (par exemple sous-schéma d'un schéma de Cohen-Macaulay). Posons
@@M1@@
et on pose aussi pour tout @@M2@@-schéma propre @@M3@@ et @@M4@@ :
@@M5@@

Si @@M6@@ est un @@M7@@-morphisme de @@M8@@-schémas propres, l'additivité des degrés de transcendance et (1.5.2) donnent immédiatement
@@M9@@
Montrons que les fonctions @@M10@@ ainsi définies satisfont aux conditions (i)-(iv) de 1.4. D'après (1.5.3), la condition (ii) est la seule non-triviale. Grâce à (iv) on peut évidemment supposer que @@M11@@. On factorise alors le morphisme structural @@M12@@ de @@M13@@ à travers son image @@M14@@, soit @@M15@@ :
@@M16@@
où @@M17@@ et @@M18@@ sont irréductibles, @@M19@@ propre et surjectif et @@M20@@ une immersion fermée. Posons @@M21@@, @@M22@@. On déduit
alors de (1.5.2) et (1.5.1) :
@@M23@@
Puisque @@M24@@ est caténaire, on a
